\documentclass[12pt,a4paper]{article}

\usepackage[brazilian]{babel}

\usepackage{amsmath}
\usepackage{amssymb}
\usepackage{enumitem}
\usepackage{etoolbox}
\usepackage{tocloft}
\usepackage{tikz}

\setlength{\cftsecnumwidth}{3.2cm}
\setlength{\cftsubsecnumwidth}{1.5cm}
\setlength{\cftsubsecnumwidth}{1.5cm}

\renewcommand{\thesection}{Capítulo \arabic{section} -}
\renewcommand{\thesubsection}{\arabic{section}.\arabic{subsection}}
\pretocmd{\subsection}{\clearpage}{}{}

\title{
    \Huge \textbf{Um Curso de Cálculo Volume 1} \\
    \vspace{0.3cm}
    \Large 6\textordfeminine{} Edição (Guidorizzi) \\
    \vspace{0.5cm}
    \large \textit{Enunciados Completos do Capítulo 3 Limite e Continuidade}
}
\author{Rogério Campos}

\begin{document}

\maketitle
\newpage

\tableofcontents
\newpage

\setcounter{section}{2}
\setcounter{subsection}{1}

\section{Limite e Continuidade}
\subsection{Introdução}
\vspace{0.5cm}
\noindent\textbf{Exercício 1.} Esboce o gráfico da função dada e, utilizando a ideia intuitiva de função contínua, determine os pontos em que a função deverá ser contínua.
\begin{enumerate}
    \renewcommand{\theenumi}{\alph{enumi}}
    \renewcommand{\labelenumi}{\theenumi)}
    \item $f(x) = 2$
    \item $f(x) = x + 1$
    \item $f(x) = x^2$
    \item $f(x) = \begin{cases} x^2 & \text{se } x \le 1 \\ 2 & \text{se } x > 1 \end{cases}$
    \item $f(x) = \begin{cases} \dfrac{1}{x^2} & \text{se } |x| \ge 1 \\ 2 & \text{se } |x| < 1 \end{cases}$
    \item $f(x) = x^2 + 2$
\end{enumerate}

\vspace{0.5cm}
\noindent\textbf{Exercício 2.} Utilizando a ideia intuitiva de limite, calcule
\begin{enumerate}
    \renewcommand{\theenumi}{\alph{enumi}}
    \renewcommand{\labelenumi}{\theenumi)}
    \item $\lim\limits_{x \to 1} (x + 2)$
    \item $\lim\limits_{x \to 1} (2x + 1)$
    \item $\lim\limits_{x \to 0} (3x + 1)$
    \item $\lim\limits_{x \to 2} (x^2 + 1)$
    \item $\lim\limits_{x \to 1} \sqrt{x}$
    \item $\lim\limits_{x \to 2} \dfrac{x^2 + x}{x + 3}$
    \item $\lim\limits_{x \to 2} \sqrt[3]{x}$
    \item $\lim\limits_{x \to 0} (\sqrt{x} + x)$
\end{enumerate}

\vspace{0.5cm}
\noindent\textbf{Exercício 3.}
Esboce o gráfico de $f(x) = \dfrac{4x^2 - 1}{2x - 1}$. Utilizando a ideia intuitiva de limite, calcule $\lim\limits_{x \to 1/2} \dfrac{4x^2 - 1}{2x - 1}$.

\vspace{0.5cm}
\noindent\textbf{Exercício 4.} Utilizando a ideia intuitiva de limite, calcule
\begin{enumerate}
    \renewcommand{\theenumi}{\alph{enumi}}
    \renewcommand{\labelenumi}{\theenumi)}
    \item $\lim\limits_{x \to 2} \dfrac{x^2 - 4}{x - 2}$
    \item $\lim\limits_{x \to 0} \dfrac{x^2 + x}{x}$
    \item $\lim\limits_{x \to 1} \dfrac{\sqrt{x} - 1}{x - 1}$
    \item $\lim\limits_{x \to 2} \dfrac{x^2 - 4x + 4}{x - 2}$
    \item $\lim\limits_{x \to -1} \dfrac{x^2 - 1}{x + 1}$
    \item $\lim\limits_{x \to 0} \operatorname{sen} x$
\end{enumerate}

\subsection{Definição de Função Contínua}
\vspace{0.5cm}
\noindent\textbf{Exercício 1.} Prove, pela definição, que a função dada é contínua no ponto dado.
\begin{enumerate}
    \renewcommand{\theenumi}{\alph{enumi}}
    \renewcommand{\labelenumi}{\theenumi)}
    \item $f(x) = 4x - 3$ em $p = 2$
    \item $f(x) = x + 1$ em $p = 2$
    \item $f(x) = -3x$ em $p = 1$
    \item $f(x) = x^3$ em $p = 2$
    \item $f(x) = x^4$ em $p = -1$
    \item $f(x) = \sqrt{x}$ em $p = 4$
    \item $f(x) = \sqrt{x}$ em $p = 0$
    \item $f(x) = \sqrt[3]{x}$ em $p = 1$
\end{enumerate}

\vspace{0.5cm}
\noindent\textbf{Exercício 2.}
Prove que $f(x) = \dfrac{1}{x}$ é contínua em todo $p \neq 0$.

\vspace{0.5cm}
\noindent\textbf{Exercício 3.}
Seja $n > 0$ um natural. Prove que $f(x) = x^n$ é contínua.

\vspace{0.5cm}
\noindent\textbf{Exercício 4.}
Prove que $f(x) = \sqrt[n]{x}$ é contínua.

\vspace{0.5cm}
\noindent\textbf{Exercício 5.}
$f(x) = \begin{cases} 2x & \text{se } x \le 1 \\ 1 & \text{se } x > 1 \end{cases}$ é contínua em $1$? Justifique.

\vspace{0.5cm}
\noindent\textbf{Exercício 6.}
Dê exemplo de uma função definida em $\mathbb{R}$ e que seja contínua em todos os pontos, exceto em $-1$, $0$, $1$.

\vspace{0.5cm}
\noindent\textbf{Exercício 7.}
Dê exemplo de uma função definida em $\mathbb{R}$ e que seja contínua em todos os pontos exceto nos inteiros.

\vspace{0.5cm}
\noindent\textbf{Exercício 8.}
Seja $f$ dada por $f(x) = \begin{cases} 1 & \text{se } x \in \mathbb{Q} \\ -1 & \text{se } x \notin \mathbb{Q} \end{cases}$. Mostre que $f$ é descontínua em todo $p$ real.

\vspace{0.5cm}
\noindent\textbf{Exercício 9.} Determine o conjunto dos pontos em que a função dada é contínua.
\begin{enumerate}
    \renewcommand{\theenumi}{\alph{enumi}}
    \renewcommand{\labelenumi}{\theenumi)}
    \item $f(x) = [\![x]\!]$ em que $[\![x]\!] = \max\{n \in \mathbb{Z} \mid n \le x\}$ (Função maior inteiro.)
    \item $f(x) = x - [\![x]\!]$
    \item $f(x) = \begin{cases} x & \text{se } x \in \mathbb{Q} \\ -x & \text{se } x \notin \mathbb{Q} \end{cases}$
    \item $f(x) = \begin{cases} x^2 - 1 & \text{se } x \in \mathbb{Q} \\ -x^2 + 1 & \text{se } x \notin \mathbb{Q} \end{cases}$
\end{enumerate}

\vspace{0.5cm}
\noindent\textbf{Exercício 10.}
Dê exemplo de uma função definida em $\mathbb{R}$ e que seja contínua apenas em $-1$, $0$, $1$.

\vspace{0.5cm}
\noindent\textbf{Exercício 11.} Determine $L$ para que a função dada seja contínua no ponto dado. Justifique.
\begin{enumerate}
    \renewcommand{\theenumi}{\alph{enumi}}
    \renewcommand{\labelenumi}{\theenumi)}
    \item $f(x) = \begin{cases} \dfrac{x^2 - 4}{x - 2} & \text{se } x \neq 2 \\ L & \text{se } x = 2 \end{cases} \quad \text{em } p = 2$
    \item $f(x) = \begin{cases} \dfrac{x^2 - x}{x} & \text{se } x \neq 0 \\ L & \text{se } x = 0 \end{cases} \quad \text{em } p = 0$
\end{enumerate}

\vspace{0.5cm}
\noindent\textbf{Exercício 12.} Dê o valor (caso exista) que a função dada deveria ter no ponto dado para ser contínua neste ponto. Justifique.
\begin{enumerate}
    \renewcommand{\theenumi}{\alph{enumi}}
    \renewcommand{\labelenumi}{\theenumi)}
    \item $g(x) = \dfrac{x^2 - 4}{x - 2} \quad \text{em } p = 2$
    \item $f(x) = \dfrac{x^2 - x}{x} \quad \text{em } p = 0$
    \item $f(x) = \dfrac{|x|}{x} \quad \text{em } p = 0$
    \item $f(x) = \begin{cases} \dfrac{x^2 - 9}{x - 3} & \text{se } x \neq 3 \\ 4 & \text{se } x = 3 \end{cases} \quad \text{em } p = 3$
    \item $g(x) = \begin{cases} x & \text{se } x < 1 \\ \dfrac{1}{x} & \text{se } x > 1 \end{cases} \quad \text{em } p = 1$
    \item $f(x) = \dfrac{|x - 2|}{x - 2} \quad \text{em } p = 2$
\end{enumerate}

\vspace{0.5cm}
\noindent\textbf{Exercício 13.}
Sabe-se que $f$ é contínua em $2$ e que $f(2) = 8$. Mostre que existe $\delta > 0$ tal que para todo $x \in D_f$
\[
2 - \delta < x < 2 + \delta \implies f(x) > 7.
\]

\vspace{0.5cm}
\noindent\textbf{Exercício 14.}
Sabe-se que $f$ é contínua em $1$ e que $f(1) = 2$. Prove que existe $r > 0$ tal que para todo $x \in D_f$
\[
1 - r < x < 1 + r \implies \dfrac{3}{2} < f(x) < \dfrac{5}{2}.
\]

\vspace{0.5cm}
\noindent\textbf{Exercício 15.}
Seja $f$ uma função definida em $\mathbb{R}$ e suponha que existe $M > 0$ tal que $|f(x) - f(p)| \le M|x - p|$ para todo $x$. Prove que $f$ é contínua em $p$.

\vspace{0.5cm}
\noindent\textbf{Exercício 16.}
Suponha que $|f(x) - f(1)| \le (x - 1)^2$ para todo $x$. Prove que $f$ é contínua em $1$.

\vspace{0.5cm}
\noindent\textbf{Exercício 17.}
Suponha que $|f(x)| \le x^2$ para todo $x$. Prove que $f$ é contínua em $0$.

\vspace{0.5cm}
\noindent\textbf{Exercício 18.}
Prove que a função $f(x) = \begin{cases} x & \text{se } x \in \mathbb{Q} \\ -x & \text{se } x \notin \mathbb{Q} \end{cases}$ é contínua em $0$.

\vspace{0.5cm}
\noindent\textbf{Exercício 19.}
Sejam $f$ e $g$ definidas em $\mathbb{R}$ e suponha que existe $M > 0$ tal que $|f(x) - f(p)| \le M|g(x) - g(p)|$ para todo $x$. Prove que se $g$ for contínua em $p$, então $f$ também será contínua em $p$.

\vspace{0.5cm}
\noindent\textbf{Exercício 20.}
Suponha $f$ definida e contínua em $\mathbb{R}$ e que $f(x) = 0$ para todo $x$ racional. Prove que $f(x) = 0$ para todo $x$ real.

\vspace{0.5cm}
\noindent\textbf{Exercício 21.}
Sejam $f$ e $g$ contínuas em $\mathbb{R}$ e tais que $f(x) = g(x)$ para todo $x$ racional. Prove que $f(x) = g(x)$ para todo $x$ real.

\vspace{0.5cm}
\noindent\textbf{Exercício 22.}
Suponha que $f$ e $g$ são contínuas em $\mathbb{R}$ e que exista $a > 0$, $a \neq 1$, tal que para todo $r$ racional, $f(r) = a^r$ e $g(r) = a^r$. Prove que $f(x) = g(x)$ em $\mathbb{R}$.

\vspace{0.5cm}
\noindent\textbf{Exercício 23.} Seja $f(x) = x + \dfrac{1}{x}$. Prove
\begin{enumerate}
    \renewcommand{\theenumi}{\alph{enumi}}
    \renewcommand{\labelenumi}{\theenumi)}
    \item $|f(x) - f(1)| \le \left(1 + \dfrac{1}{x}\right)|x - 1|$ para $x > 0$
    \item $|f(x) - f(1)| \le 3|x - 1|$ para $x > \dfrac{1}{2}$
    \item $f$ é contínua em $p = 1$
\end{enumerate}

\vspace{0.5cm}
\noindent\textbf{Exercício 24.} Seja $f(x) = x^3 + x$. Prove que
\begin{enumerate}
    \renewcommand{\theenumi}{\alph{enumi}}
    \renewcommand{\labelenumi}{\theenumi)}
    \item $|f(x) - f(2)| \le 20|x - 2|$ para $0 \le x \le 3$
    \item $f$ é contínua em $2$
\end{enumerate}

\vspace{0.5cm}
\noindent\textbf{Exercício 25.}
Prove que $f(x) = x + \dfrac{1}{x^2}$ é contínua em $1$.

\vspace{0.5cm}
\noindent\textbf{Exercício 26.}
Prove que $f(x) = x + \dfrac{1}{x}$ é contínua em todo $p > 0$.

\vspace{0.5cm}
\noindent\textbf{Exercício 27.} Sejam $f(x) = x^3$ e $p \neq 0$.
\begin{enumerate}
    \renewcommand{\theenumi}{\alph{enumi}}
    \renewcommand{\labelenumi}{\theenumi)}
    \item Verifique que $|x^3 - p^3| \le 7\,p^2|x - p|$ para $|x| \le 2|p|$
    \item Conclua de (a) que $f$ é contínua em $p$
\end{enumerate}

\subsection{Definição de Limite}
\vspace{0.5cm}
\noindent\textbf{Exercício 1.} Calcule e justifique.
\begin{enumerate}
    \renewcommand{\theenumi}{\alph{enumi}}
    \renewcommand{\labelenumi}{\theenumi)}
    \item $\lim\limits_{x \to 2} x^2$
    \item $\lim\limits_{x \to 1} (3x + 1)$
    \item $\lim\limits_{x \to -2} (4x + 1)$
    \item $\lim\limits_{x \to 10} 5$
    \item $\lim\limits_{x \to -9} 50$
    \item $\lim\limits_{x \to -1} (-x^2 - 2x + 3)$
    \item $\lim\limits_{x \to 4} \sqrt{x}$
    \item $\lim\limits_{x \to -3} \sqrt[3]{x}$
    \item $\lim\limits_{x \to -8} \sqrt{5}$
    \item $\lim\limits_{x \to 3} \dfrac{x^2 - 9}{x - 3}$
    \item $\lim\limits_{x \to 3} \dfrac{x^2 - 9}{x + 3}$
    \item $\lim\limits_{x \to -1} \dfrac{x^2 - 9}{x - 3}$
    \item $\lim\limits_{x \to \frac{1}{2}} \dfrac{4x^2 - 1}{2x - 1}$
    \item $\lim\limits_{x \to 1} \dfrac{\sqrt{x} - 1}{x - 1}$
    \item $\lim\limits_{x \to -\frac{1}{3}} \dfrac{9x^2 - 1}{3x + 1}$
    \item $\lim\limits_{x \to 3} \dfrac{\sqrt{x} - \sqrt{3}}{x - 3}$
    \item $\lim\limits_{x \to 3} \dfrac{\sqrt[3]{x} - \sqrt[3]{3}}{x - 3}$
    \item $\lim\limits_{x \to 2} \dfrac{\sqrt[4]{x} - \sqrt[4]{2}}{x - 2}$
    \item $\lim\limits_{x \to 0} \dfrac{x^2 + 3x - 1}{x^2 + 2}$
    \item $\lim\limits_{x \to 1} \dfrac{\sqrt{x} - 1}{\sqrt{2x + 3} - \sqrt{5}}$
\end{enumerate}

\vspace{0.5cm}
\noindent\textbf{Exercício 2.} Determine $L$ para que a função dada seja contínua no ponto dado. Justifique.
\begin{enumerate}
    \renewcommand{\theenumi}{\alph{enumi}}
    \renewcommand{\labelenumi}{\theenumi)}
    \item $f(x) = \begin{cases} \dfrac{x^3 - 8}{x - 2} & \text{se } x \neq 2 \\ L & \text{se } x = 2 \end{cases} \quad \text{em } p = 2$
    \item $f(x) = \begin{cases} \dfrac{\sqrt{x} - \sqrt{3}}{x - 3} & \text{se } x \neq 3 \\ L & \text{se } x = 3 \end{cases} \quad \text{em } p = 3$
    \item $f(x) = \begin{cases} \dfrac{\sqrt{x} - \sqrt{5}}{\sqrt{x + 5} - \sqrt{10}} & \text{se } x \neq 5 \\ L & \text{se } x = 5 \end{cases} \quad \text{em } p = 5$
\end{enumerate}

\vspace{0.5cm}
\noindent\textbf{Exercício 3.}
$f(x) = \begin{cases} \dfrac{x^2 + x}{x + 1} & \text{se } x \neq -1 \\ 2 & \text{se } x = -1 \end{cases}$ é contínua em $-1$? E em $0$? Por quê?

\vspace{0.5cm}
\noindent\textbf{Exercício 4.} Calcule $\lim\limits_{h \to 0} \dfrac{f(x + h) - f(x)}{h}$ sendo $f$ dada por
\begin{enumerate}
    \renewcommand{\theenumi}{\alph{enumi}}
    \renewcommand{\labelenumi}{\theenumi)}
    \item $f(x) = x^2$
    \item $f(x) = 2x^2 + x$
    \item $f(x) = 5$
    \item $f(x) = -x^3 + 2x$
    \item $f(x) = \dfrac{1}{x}$
    \item $f(x) = 3x + 1$
\end{enumerate}

\vspace{0.5cm}
\noindent\textbf{Exercício 5.} Calcule.
\begin{enumerate}
    \renewcommand{\theenumi}{\alph{enumi}}
    \renewcommand{\labelenumi}{\theenumi)}
    \item $\lim\limits_{x \to -1} \dfrac{x^3 + 1}{x^2 - 1}$
    \item $\lim\limits_{x \to 0} \dfrac{x^3 + x^2}{3x^3 + x^4 + x}$
    \item $\lim\limits_{h \to 0} (x^2 + 3xh)$
    \item $\lim\limits_{h \to 0} \dfrac{(x + h)^3 - x^3}{h}$
    \item $\lim\limits_{x \to 3} \dfrac{x^2 - 9}{x^2 + 9}$
    \item $\lim\limits_{x \to p} \dfrac{\sqrt[3]{x} - \sqrt[3]{p}}{x - p} \quad (p \neq 0)$
    \item $\lim\limits_{x \to p} \dfrac{\sqrt[4]{x} - \sqrt[4]{p}}{x - p} \quad (p \neq 0)$
    \item $\lim\limits_{x \to 2} \dfrac{x^3 - 5x^2 + 8x - 4}{x^4 - 5x - 6}$
    \item $\lim\limits_{x \to 1} \dfrac{x^3 - 1}{x^4 + 3x - 4}$
    \item $\lim\limits_{x \to 7} \dfrac{\sqrt{x} - \sqrt{7}}{\sqrt{x + 7} - \sqrt{14}}$
    \item $\lim\limits_{x \to p} \dfrac{x^3 - p^3}{x - p}$
    \item $\lim\limits_{x \to p} \dfrac{x^4 - p^4}{x - p}$
    \item $\lim\limits_{x \to p} \dfrac{x^n - p^n}{x - p} \quad (n > 0 \text{ natural})$
    \item $\lim\limits_{x \to p} \dfrac{\sqrt[n]{x} - \sqrt[n]{p}}{x - p}$
    \item $\lim\limits_{x \to 2} \dfrac{\dfrac{1}{x} - \dfrac{1}{2}}{x - 2}$
    \item $\lim\limits_{x \to p} \dfrac{f(x) - f(p)}{x - p}$ em que $f(x) = \dfrac{1}{x}$
    \item $\lim\limits_{x \to p} \dfrac{g(x) - g(p)}{x - p}$ em que $g(x) = \dfrac{1}{x^2}$
    \item $\lim\limits_{h \to 0} \dfrac{f(x + h) - f(x)}{h}$ em que $f(x) = x^2 - 3x$
\end{enumerate}

\vspace{0.5cm}
\noindent\textbf{Exercício 6.}
Prove que existe $\delta > 0$ tal que
\[
1 - \delta < x < 1 + \delta \implies 2 - \dfrac{1}{3} < x^2 + x < 2 + \dfrac{1}{3}.
\]

\vspace{0.5cm}
\noindent\textbf{Exercício 7.}
Prove que existe $\delta > 0$ tal que
\[
1 - \delta < x < 1 + \delta \implies 2 - \dfrac{1}{2} < \dfrac{x^5 + 3x}{x^2 + 1} < 2 + \dfrac{1}{2}.
\]

\vspace{0.5cm}
\noindent\textbf{Exercício 8.}
Sejam $f$ e $g$ definidas em $\mathbb{R}$ com $g(x) \neq 0$ para todo $x$. Suponha que $\lim\limits_{x \to p} \dfrac{f(x)}{g(x)} = 0$. Prove que existe $\delta > 0$ tal que
\[
0 < |x - p| < \delta \implies |f(x)| < |g(x)|.
\]

\vspace{0.5cm}
\noindent\textbf{Exercício 9.}
Suponha que $\lim\limits_{x \to p} f(x) = L$. Prove que existem $r > 0$, $\alpha$ e $\beta$ tais que, para todo $x \in D_f$,
\[
0 < |x - p| < r \implies \alpha < f(x) < \beta.
\]
Interprete graficamente.

\vspace{0.5cm}
\noindent\textbf{Exercício 10.}
Suponha que $\lim\limits_{x \to p} f(x) = L$. Prove que existem $r > 0$ e $M > 0$ tais que, para todo $x \in D_f$,
\[
0 < |x - p| < r \implies |f(x)| \le M.
\]

\vspace{0.5cm}
\noindent\textbf{Exercício 11.}
Prove: $\lim\limits_{x \to p} f(x) = L \iff \lim\limits_{x \to p} [f(x) - L] = 0$.

\vspace{0.5cm}
\noindent\textbf{Exercício 12.}
Prove: $\lim\limits_{x \to p} f(x) = L \iff \lim\limits_{x \to p} |f(x) - L| = 0$.

\vspace{0.5cm}
\noindent\textbf{Exercício 13.}
Prove: $\lim\limits_{x \to p} \dfrac{f(x)}{x - p} = 0 \iff \lim\limits_{x \to p} \dfrac{f(x)}{|x - p|} = 0$.

\vspace{0.5cm}
\noindent\textbf{Exercício 14.}
Suponha que existe $r > 0$ tal que $f(x) \ge 0$ para $0 < |x - p| < r$ e que $\lim\limits_{x \to p} f(x) = L$. Prove que $L \ge 0$. (Sugestão: Suponha $L < 0$ e use a conservação do sinal.)

\vspace{0.5cm}
\noindent\textbf{Exercício 15.}
Suponha $f$ contínua em $\mathbb{R}$ e $f(x) \ge 0$ para todo $x$ racional. Prove que $f(x) \ge 0$ para todo $x$.

\subsection{Limites Laterais}
\vspace{0.5cm}
\noindent\textbf{Exercício 1.} Calcule, caso exista. Se não existir, justifique.
\begin{enumerate}
    \renewcommand{\theenumi}{\alph{enumi}}
    \renewcommand{\labelenumi}{\theenumi)}
    \item $\lim\limits_{x \to 1^+} \dfrac{|x - 1|}{x - 1}$
    \item $\lim\limits_{x \to 1^-} \dfrac{|x - 1|}{x - 1}$
    \item $\lim\limits_{x \to 1^+} \dfrac{f(x) - f(1)}{x - 1}$ em que $f(x) = \begin{cases} x + 1 & \text{se } x \ge 1 \\ 2x & \text{se } x < 1 \end{cases}$
    \item $\lim\limits_{x \to 0} \sqrt{x}$
    \item $\lim\limits_{x \to 1} \dfrac{|x - 1|}{x - 1}$
    \item $\lim\limits_{x \to 1^-} \dfrac{f(x) - f(1)}{x - 1}$ em que $f(x) = \begin{cases} x + 1 & \text{se } x \ge 1 \\ 2x & \text{se } x < 1 \end{cases}$
    \item $\lim\limits_{x \to 2^+} \dfrac{x^2 - 2x + 1}{x - 1}$
    \item $\lim\limits_{x \to 3} \dfrac{|x - 1|}{x - 1}$
    \item $\lim\limits_{x \to 1} \dfrac{f(x) - f(1)}{x - 1}$ em que $f(x) = \begin{cases} x^2 & \text{se } x \le 1 \\ 2x - 1 & \text{se } x > 1 \end{cases}$
    \item $\lim\limits_{x \to 2^-} \dfrac{g(x) - g(2)}{x - 2}$ em que $g(x) = \begin{cases} x & \text{se } x \ge 2 \\ \dfrac{x^2}{2} & \text{se } x < 2 \end{cases}$
    \item $\lim\limits_{x \to 2^+} \dfrac{g(x) - g(2)}{x - 2}$ sendo $g$ a função do item (j)
    \item $\lim\limits_{x \to 2} \dfrac{g(x) - g(2)}{x - 2}$ em que $g$ é a função do item (j)
\end{enumerate}

\vspace{0.5cm}
\noindent\textbf{Exercício 2.}
A afirmação ``$\lim\limits_{x \to p^+} f(x) = \lim\limits_{x \to p^-} f(x) \implies f \text{ contínua em } p$'' é falsa ou verdadeira? Justifique.

\vspace{0.5cm}
\noindent\textbf{Exercício 3.}
Dada a função $f(x) = \dfrac{x^2 - 3x + 2}{x - 1}$, verifique que $\lim\limits_{x \to 1^+} f(x) = \lim\limits_{x \to 1^-} f(x)$. Pergunta-se: $f$ é contínua em $1$? Por quê?

\vspace{0.5cm}
\noindent\textbf{Exercício 4.}
Dê exemplo de uma função definida em $\mathbb{R}$, que não seja contínua em $2$, mas que $\lim\limits_{x \to 2^+} f(x) = \lim\limits_{x \to 2^-} f(x)$.

\vspace{0.5cm}
\noindent\textbf{Exercício 5.}
Suponha que exista $r > 0$ tal que $f(x) \ge 0$ para $p < x < p + r$. Prove que $\lim\limits_{x \to p^+} f(x) \ge 0$ desde que o limite exista.

\vspace{0.5cm}
\noindent\textbf{Exercício 6.}
Sejam $f$ uma função definida em um intervalo aberto $I$ e $p \in I$. Suponha que $f(x) \le f(p)$ para todo $x \in I$. Prove que $\lim\limits_{x \to p} \dfrac{f(x) - f(p)}{x - p} = 0$ desde que o limite exista.
(Sugestão: Estude os sinais de $\lim\limits_{x \to p^+} \dfrac{f(x) - f(p)}{x - 1}$ e de $\lim\limits_{x \to p^-} \dfrac{f(x) - f(p)}{x - 1}$.)

\subsection{Limite de Função Composta}
\vspace{0.5cm}
\noindent\textbf{Exercício 1.} Calcule
\begin{enumerate}
    \renewcommand{\theenumi}{\alph{enumi}}
    \renewcommand{\labelenumi}{\theenumi)}
    \item $\lim\limits_{x \to -1} \sqrt[3]{\dfrac{x^3 + 1}{x + 1}}$
    \item $\lim\limits_{x \to 1} \dfrac{\sqrt{x^2 + 3} - 2}{x^2 - 1}$
    \item $\lim\limits_{x \to 1} \dfrac{\sqrt[3]{x + 7} - 2}{x - 1}$
    \item $\lim\limits_{x \to 1} \dfrac{\sqrt[3]{3x + 5} - 2}{x^2 - 1}$
\end{enumerate}

\vspace{0.5cm}
\noindent\textbf{Exercício 2.} Seja $f$ definida em $\mathbb{R}$. Suponha que $\lim\limits_{x \to 0} \dfrac{f(x)}{x} = 1$. Calcule
\begin{enumerate}
    \renewcommand{\theenumi}{\alph{enumi}}
    \renewcommand{\labelenumi}{\theenumi)}
    \item $\lim\limits_{x \to 0} \dfrac{f(3x)}{x}$
    \item $\lim\limits_{x \to 0} \dfrac{f(x^2)}{x}$
    \item $\lim\limits_{x \to 1} \dfrac{f(x^2 - 1)}{x - 1}$
    \item $\lim\limits_{x \to 0} \dfrac{f(7x)}{3x}$
\end{enumerate}

\vspace{0.5cm}
\noindent\textbf{Exercício 3.} Seja $f$ definida em $\mathbb{R}$ e seja $p$ um real dado. Suponha que $\lim\limits_{x \to p} \dfrac{f(x) - f(p)}{x - p} = L$. Calcule
\begin{enumerate}
    \renewcommand{\theenumi}{\alph{enumi}}
    \renewcommand{\labelenumi}{\theenumi)}
    \item $\lim\limits_{h \to 0} \dfrac{f(p + h) - f(p)}{h}$
    \item $\lim\limits_{h \to 0} \dfrac{f(p + 3h) - f(p)}{h}$
    \item $\lim\limits_{h \to 0} \dfrac{f(p + h) - f(p - h)}{h}$
    \item $\lim\limits_{h \to 0} \dfrac{f(p - h) - f(p)}{h}$
\end{enumerate}

\subsection{Teorema do Confronto}
\vspace{0.5cm}
\noindent\textbf{Exercício 1.}
Seja $f$ uma função definida em $\mathbb{R}$ tal que, para todo $x \neq 1$, $-x^2 + 3x \le f(x) < \dfrac{x^2 - 1}{x - 1}$. Calcule $\lim\limits_{x \to 1} f(x)$ e justifique.

\vspace{0.5cm}
\noindent\textbf{Exercício 2.}
Seja $f$ definida em $\mathbb{R}$ e tal que, para todo $x$, $|f(x) - 3| \le 2|x - 1|$. Calcule $\lim\limits_{x \to 1} f(x)$ e justifique.

\vspace{0.5cm}
\noindent\textbf{Exercício 3.}
Suponha que, para todo $x$, $|g(x)| \le x^4$. Calcule $\lim\limits_{x \to 0} \dfrac{g(x)}{x}$.

\vspace{0.5cm}
\noindent\textbf{Exercício 4.}
\begin{enumerate}
    \renewcommand{\theenumi}{\alph{enumi}}
    \renewcommand{\labelenumi}{\theenumi)}
    \item Verifique que $\lim\limits_{x \to 0} \operatorname{sen}\dfrac{1}{x}$ não existe.
    \item Calcule, caso exista, $\lim\limits_{x \to 0} x\operatorname{sen}\dfrac{1}{x}$. (Justifique.)
\end{enumerate}

\vspace{0.5cm}
\noindent\textbf{Exercício 5.} Calcule, caso exista, $\lim\limits_{x \to 0} \dfrac{f(x) - f(0)}{x - 0}$ em que $f$ é dada por
\begin{enumerate}
    \renewcommand{\theenumi}{\alph{enumi}}
    \renewcommand{\labelenumi}{\theenumi)}
    \item $f(x) = \begin{cases} x^2 \operatorname{sen}\dfrac{1}{x} & \text{se } x \neq 0 \\ 0 & \text{se } x = 0 \end{cases}$
    \item $f(x) = \begin{cases} x \operatorname{sen}\dfrac{1}{x} & \text{se } x \neq 0 \\ 0 & \text{se } x = 0 \end{cases}$
\end{enumerate}

\vspace{0.5cm}
\noindent\textbf{Exercício 6.} Sejam $f$ e $g$ duas funções definidas em $\mathbb{R}$ e tais que, para todo $x$, $[g(x)]^4 + [f(x)]^4 = 4$. Calcule e justifique.
\begin{enumerate}
    \renewcommand{\theenumi}{\alph{enumi}}
    \renewcommand{\labelenumi}{\theenumi)}
    \item $\lim\limits_{x \to 0} x^3 g(x)$
    \item $\lim\limits_{x \to 3} f(x)\sqrt[3]{x^2 - 9}$
\end{enumerate}

\vspace{0.5cm}
\noindent\textbf{Exercício 7.} Seja $f$ definida em $\mathbb{R}$ e suponha que existe $M > 0$ tal que, para todo $x$, $|f(x) - f(p)| \le M|x - p|^2$.
\begin{enumerate}
    \renewcommand{\theenumi}{\alph{enumi}}
    \renewcommand{\labelenumi}{\theenumi)}
    \item Mostre que $f$ é contínua em $p$.
    \item Calcule, caso exista, $\lim\limits_{x \to p} \dfrac{f(x) - f(p)}{x - p}$.
\end{enumerate}

\vspace{0.5cm}
\noindent\textbf{Exercício 8.}
Sejam $a$, $b$, $c$ reais fixos e suponha que, para todo $x$, $|a + bx + cx^2| \le |x|^3$. Prove que $a = b = c = 0$.

\vspace{0.5cm}
\noindent\textbf{Exercício 9.}
Prove: $\lim\limits_{x \to p} f(x) = L \implies \lim\limits_{x \to p} |f(x)| = |L|$.\\
(Sugestão: verifique que $\big||f(x)| - |L|\big| \le |f(x) - L|$ e aplique o teorema do confronto.)

\vspace{0.5cm}
\noindent\textbf{Exercício 10.}
A afirmação
\[
\text{``}\lim_{x \to p} |f(x)| = |L| \implies \lim_{x \to p} f(x) = L\text{''}
\]
é falsa ou verdadeira? Por quê?

\vspace{0.5cm}
\noindent\textbf{Exercício 11.}
Dê exemplo de uma função $f$ tal que $\lim\limits_{x \to p} |f(x)|$ existe, mas $\lim\limits_{x \to p} f(x)$ não exista.

\vspace{0.5cm}
\noindent\textbf{Exercício 12.}
Prove: $\lim\limits_{h \to 0} \dfrac{f(h)}{h} = 0 \iff \lim\limits_{h \to 0} \dfrac{f(h)}{|h|} = 0$.

\subsection{Continuidade das Funções Trigonométricas}
Esta subseção não possui exercícios

\subsection{O Limite Fundamental $\lim_{x \to 0} \frac{sen(x)}{x}$}
\vspace{0.5cm}
\noindent\textbf{Exercício 1.} Calcule.
\begin{enumerate}
    \renewcommand{\theenumi}{\alph{enumi}}
    \renewcommand{\labelenumi}{\theenumi)}
    \item $\lim\limits_{x \to 0} \dfrac{\operatorname{tg} x}{x}$
    \item $\lim\limits_{x \to 0} \dfrac{x}{\operatorname{sen} x}$
    \item $\lim\limits_{x \to 0} \dfrac{\operatorname{sen} 3x}{x}$
    \item $\lim\limits_{x \to \pi} \dfrac{\operatorname{sen} x}{x - \pi}$
    \item $\lim\limits_{x \to 0} \dfrac{x^2}{\operatorname{sen} x}$
    \item $\lim\limits_{x \to 0} \dfrac{3x^2}{\operatorname{tg} x \operatorname{sen} x}$
    \item $\lim\limits_{x \to 0} \dfrac{\operatorname{tg} 3x}{\operatorname{sen} 4x}$
    \item $\lim\limits_{x \to 0} \dfrac{1 - \cos x}{x}$
    \item $\lim\limits_{x \to \frac{\pi}{2}} \dfrac{1 - \operatorname{sen} x}{2x - \pi}$
    \item $\lim\limits_{x \to 0} x \operatorname{sen}\dfrac{1}{x}$
    \item $\lim\limits_{x \to p} \dfrac{\operatorname{tg}(x - p)}{x^2 - p^2}$, $p \neq 0$
    \item $\lim\limits_{x \to p} \dfrac{\operatorname{sen}(x^2 - p^2)}{x - p}$
    \item $\lim\limits_{x \to 0} \dfrac{\operatorname{sen}\!\left(x^2 + \dfrac{1}{x}\right) - \operatorname{sen}\dfrac{1}{x}}{x}$
    \item $\lim\limits_{x \to 0} \dfrac{x + \operatorname{sen} x}{x^2 - \operatorname{sen} x}$
    \item $\lim\limits_{x \to 0} \dfrac{x - \operatorname{tg} x}{x + \operatorname{tg} x}$
    \item $\lim\limits_{x \to 1} \dfrac{\operatorname{sen}\pi x}{x - 1}$
\end{enumerate}

\vspace{0.5cm}
\noindent\textbf{Exercício 2.}
\begin{enumerate}
    \renewcommand{\theenumi}{\alph{enumi}}
    \renewcommand{\labelenumi}{\theenumi)}
    \item Prove que existe $r > 0$ tal que
    \[
    \cos x - 1 < \dfrac{\operatorname{sen} x}{x} - 1 < 0
    \]
    para $0 < |x| < r$.
    \item Calcule $\lim\limits_{x \to 0} \dfrac{x - \operatorname{sen} x}{x^2}$.
\end{enumerate}

\vspace{0.5cm}
\noindent\textbf{Exercício 3.} Calcule.
\begin{enumerate}
    \renewcommand{\theenumi}{\alph{enumi}}
    \renewcommand{\labelenumi}{\theenumi)}
    \item $\lim\limits_{x \to p} \dfrac{\operatorname{sen} x - \operatorname{sen} p}{x - p}$
    \item $\lim\limits_{x \to p} \dfrac{\cos x - \cos p}{x - p}$
    \item $\lim\limits_{x \to p} \dfrac{\operatorname{tg} x - \operatorname{tg} p}{x - p}$
    \item $\lim\limits_{x \to p} \dfrac{\sec x - \sec p}{x - p}$
\end{enumerate}

\subsection{Propriedades Operatórias}
Esta subseção não possui exercícios

\end{document}
